\documentclass[fleqn]{article}
\oddsidemargin 0.0in
\textwidth 6.0in
\thispagestyle{empty}
\usepackage{import}
\usepackage{amsmath}
\usepackage{graphicx}
\usepackage{flexisym}
\usepackage{calligra}
\usepackage{amssymb}
\usepackage{bigints} 
\usepackage[english]{babel}
\usepackage[utf8x]{inputenc}
\usepackage{float}
\usepackage[colorinlistoftodos]{todonotes}


\DeclareMathAlphabet{\mathcalligra}{T1}{calligra}{m}{n}
\DeclareFontShape{T1}{calligra}{m}{n}{<->s*[2.2]callig15}{}
\newcommand{\scriptr}{\mathcalligra{r}\,}
\newcommand{\boldscriptr}{\pmb{\mathcalligra{r}}\,}

\definecolor{hwColor}{HTML}{442020}

\begin{document}

  \begin{titlepage}

    \newcommand{\HRule}{\rule{\linewidth}{0.5mm}}

    \center

    \begin{center}
      \includegraphics[height=11cm, width=11cm]{asu.png}
    \end{center}

    \vline

    \textsc{\LARGE Statistical/Thermal Physics}\\[1.5cm]

    \HRule \\[0.5cm]
    { \huge \bfseries Homework 5}\\[0.4cm] 
    \HRule \\[1.0cm]

    \textbf{Behnam Amiri}

    \bigbreak

    \textbf{Prof: Michael Treacy}

    \bigbreak

    \textbf{{\large \today}\\[2cm]}

    \vfill

  \end{titlepage}

  \begin{enumerate}
    \item \textbf{2.27} Rather than insisting that all the molecules be in the left half of a container, suppose we only require 
    that they be in the leftmost $99\%$ (leaving the remaining $1\%$ completely empty). What is the probability of finding such an
    arrangement if there are 100 molecules in the container? What if there are $10,000$ molecules? What if there are $10^{23}$?

      % \textcolor{hwColor}{
      %   \\
      % }

    \item \textbf{2.28} How many possible arrangement are there for a deck of 52 playing cards?..... Is this entrophy significant
    compared to the entrophy associated with arranging thermal energy amoung the molecules in the cards?

      % \textcolor{hwColor}{
      %   \\
      % }

    \item \textbf{2.29} Consider a system of two Einstein  solids, with $N_A=300, ~ N_B=200$, and $q_{total}=100$ (as discussed in Section
    2.3). Compute the entrophy of the most likely macrostate...

      % \textcolor{hwColor}{
      %   \\
      % }

    \item \textbf{2.30} Consider again the system of two large, identical Einstein solids treated in Problem 2.22.
    \begin{enumerate}
      \item For the case $N=10^{23}$, compute the entrophy of this system...

        % \textcolor{hwColor}{
        %   \\
        % }

      \item Compute the entrophy again, assuming that the system is in its most likely macrostate...

        % \textcolor{hwColor}{
        %   \\
        % }

      \item Is the issue of time scales really relevant to the entrophy of this system?

        % \textcolor{hwColor}{
        %   \\
        % }
      
      \item Suppose that, at a moment when the system is near its most likely macrostate, you suddenly insert a partition between 
      the solids so that they can no longer...

        % \textcolor{hwColor}{
        %   \\
        % }

    \end{enumerate}


    \item \textbf{2.33} Use the Sackur-Tetrode equation to calculate the entrophy of a mole of argon gas at room 
    temperature and atmospheric pressure. Why is the entrophy greater than that of a mole of helium under the same conditions?

      % \textcolor{hwColor}{
      %   \\
      % }

    \item \textbf{2.34} Show that during the quasistatic isothermal expansion of a monatomic ideal gas, the change in entrophy is
    related to the heat input $Q$ by the simple formula
    $$
      \Delta S=\dfrac{Q}{T}
    $$
    In the following chapter, I will prove that this formula is valid for any quasistatic process. Show, however, that it is not valid
    for the free expansion process described above.

      % \textcolor{hwColor}{
      %   \\
      % }

    \item \textbf{2.37} Using the same method as in the text....
    $$  
      \Delta S_{mixing}=-Nk \left[x ~ ln\bigg( x \bigg)+\bigg( 1-x \bigg) ~ ln\bigg( 1-x \bigg)\right]
    $$
    Check that this expression reduces to the one given in the text when $x=\dfrac{1}{2}$.

      % \textcolor{hwColor}{
      %   \\
      % }

    \item \textbf{2.39} Compute the entrophy of a mole of helium at room temperature and atmospheric pressure, pretending that all 
    the atoms are distinguishable. Compare to the actual entrophy, for indistinguishable atoms, computed in the text.

      % \textcolor{hwColor}{
      %   \\
      % }

  \end{enumerate}

\end{document}
